\documentclass[11pt,a4paper]{article}

% ============================================================
%  Packages
% ============================================================
\usepackage[utf8]{inputenc}
\usepackage[T1]{fontenc}
\usepackage[english]{babel}
\usepackage{lmodern}
\usepackage{microtype}
\usepackage[a4paper,margin=2.5cm]{geometry}
\usepackage{amsmath,amssymb,amsthm,mathtools}
\usepackage{bm}
\usepackage{graphicx}
\usepackage{booktabs}
\usepackage{array}
\usepackage{longtable}
\usepackage{multirow}
\usepackage{siunitx}
\usepackage{enumitem}
\usepackage{hyperref}
\usepackage[numbers,sort&compress]{natbib}
\bibliographystyle{plainnat}
\usepackage{xcolor}
\usepackage{caption}
\captionsetup{font=small,labelfont=bf}

\hypersetup{
  colorlinks = true,
  linkcolor  = blue!50!black,
  citecolor  = green!40!black,
  urlcolor   = blue!70!black,
  pdftitle   = {The Complete Atom from First Principles},
  pdfauthor  = {Kundai Farai Sachikonye}
}

% ============================================================
%  Theorem environments
% ============================================================
\newtheorem{theorem}{Theorem}[section]
\newtheorem{lemma}[theorem]{Lemma}
\newtheorem{proposition}[theorem]{Proposition}
\newtheorem{corollary}[theorem]{Corollary}

\theoremstyle{definition}
\newtheorem{definition}[theorem]{Definition}
\newtheorem{axiom}[theorem]{Axiom}
\newtheorem{example}[theorem]{Example}

\theoremstyle{remark}
\newtheorem{remark}[theorem]{Remark}

% ============================================================
%  Custom commands
% ============================================================
\newcommand{\Sk}{\mathcal{S}_{\mathrm{cat}}}
\newcommand{\St}{\mathcal{S}_{\mathrm{osc}}}
\newcommand{\Se}{\mathcal{S}_{\mathrm{part}}}
\newcommand{\Sspace}{\mathcal{S}}
\newcommand{\kB}{k_{\scriptscriptstyle B}}
\newcommand{\Depth}{\mathcal{M}}
\newcommand{\Real}{\mathbb{R}}
\newcommand{\Natural}{\mathbb{N}}
\newcommand{\Integer}{\mathbb{Z}}
\newcommand{\Hilbert}{\mathcal{H}}
\newcommand{\Tp}{T_{\scriptscriptstyle P}}
\newcommand{\Lp}{\ell_{\scriptscriptstyle P}}

% ============================================================
%  Title
% ============================================================
\title{\textbf{The Complete Atom from First Principles:\\[2pt]
Sub-Nuclear Virtual Substates, Nuclear Partition Structure,\\ and Electronic Shells Derived from Bounded Phase Space Geometry\\ and Validated Against Foundational Experiments}}

\author{Kundai Farai Sachikonye\\
AIMe Registry for Artificial Intelligence\\
\texttt{kundai.sachikonye@bitspark.com}}

\date{\today}

\begin{document}
\maketitle

% ============================================================
%  Abstract
% ============================================================
\begin{abstract}
We derive the complete atom --- from the sub-nucleon substrate through the nucleon, the nucleus, the nucleus--electron system, and the observable atomic volume --- from a single axiom: physical systems occupy bounded regions of phase space with finite Liouville measure, and these bounded regions admit hierarchical partitioning into distinguishable subregions. From this axiom, combined with the Loschmidt principle that identity is relational, we derive without postulate: (i)~oscillatory dynamics as the only mode of bounded existence; (ii)~the Triple Equivalence Theorem showing that oscillatory, categorical, and partition descriptions are mathematically identical with common state function $\mathcal{S} = \kB M \ln n$; (iii)~the partition depth field $\Depth(\mathbf{x})$ whose gradient replaces all forces; (iv)~partition coordinates $(n,l,m,s)$ with capacity $C(n)=2n^2$; (v)~a continuous embedding theorem requiring a $[0,1]^3$ Riemannian manifold for backward navigation and introducing sub-coordinate virtual substates; and (vi)~a charge emergence theorem under which charge exists only at partition boundaries. Applied to the atom, these results yield: sub-nucleon constituents as structurally unobservable virtual substates (reproducing confinement without naming it); the nucleon as a bounded system of finite radius $\approx 0.84$~fm with exponential charge profile matching Hofstadter's form factors; the nucleus as a close-packed partition aggregate with $R = r_0 A^{1/3}$, $r_0 \approx 1.2$~fm; binding energies from the Composition Theorem $(\Delta\Depth)$ reproducing $^2$H (2.22~MeV), $^4$He (28.3~MeV), $^{16}$O (127.6~MeV), $^{56}$Fe (492.3~MeV) and $^{238}$U (1801.7~MeV); magic numbers 2, 8, 20, 28, 50, 82, 126 from nuclear partition shells with geometric spin--orbit coupling; the electron as a partition depth minimum around the $+Ze$ boundary produced by measurement; a Coulomb-form gradient $-\nabla\Depth \propto Ze^2/(4\pi\epsilon_0 r^2)$ without invoking a Coulomb force; effective charge screening matching Slater's rules; and the geometric meaning of atomic volume as the region the electron is ``not at''. A major implication is that the intra-nuclear Coulomb repulsion paradox is dissolved, not solved: unpartitioned nucleons inside the nucleus have no individual charge, so no repulsion exists to require a strong force. We revisit twelve foundational experiments --- Rutherford (1911), Hofstadter (1956), Moseley (1913), Millikan (1913), Einstein photoelectric (1905), Compton (1923), Stern--Gerlach (1922), Zeeman (1897), Franck--Hertz (1914), Davisson--Germer (1927), the binding energy per nucleon curve, and deep inelastic scattering --- and show that each is reproduced quantitatively while being interpreted as direct evidence for the partition structure. The framework carries zero adjustable parameters.
\end{abstract}

\vspace{0.5em}
\textbf{Keywords:} bounded phase space; partition coordinates; nuclear structure; binding energy; magic numbers; form factors; atomic volume; Rutherford scattering; foundational experiments; charge emergence.

\tableofcontents

% ============================================================
%                        PART I
% ============================================================
\part{Foundations}

\section{Introduction}
\label{sec:intro}

The modern account of the atom rests on an uncomfortable collection of independent postulates. The nucleus is held together by a postulated ``strong force'' invoked to overcome an equally postulated intra-nuclear Coulomb repulsion. The Coulomb interaction is itself a long-range $1/r^2$ field said to fill the atomic volume, yet in Rutherford's own scattering experiment \citep{rutherford1911} the majority of $\alpha$ particles traverse the atom without the slightest detectable deflection --- behaviour inconsistent with a genuine force field continuously pervading the atomic interior. Electrons are said to occupy shells derived from the Schr\"odinger equation \citep{schrodinger1926}, yet the 4-tuple $(n,l,m,s)$ itself is not derived from the equation; it is imposed as an interpretation of the spectrum. Wave--particle duality stands beside quantisation stands beside the Pauli principle \citep{pauli1925} as three independent principles, each required, none derived from the others. Mass is identified with rest energy via $m = E/c^2$ but the origin of $c$ is not itself explained.

We show in this paper that every one of these ingredients is unnecessary. A single axiom --- that physical systems occupy bounded regions of phase space with finite Liouville measure admitting hierarchical partition --- together with the Loschmidt principle that identity is relational, suffices to derive the complete atom quantitatively: the sub-nucleon substrate, the nucleon itself, the nucleus with its binding energies and magic numbers, the electron and its shells around the nucleus, and the geometric meaning of atomic volume. No force is postulated. No charge is postulated. No wave--particle duality is postulated. Each of these emerges as a theorem in the framework.

The paper is organised in nine parts. Part~I states and proves the foundations: the axiom, the Loschmidt principle, the Triple Equivalence Theorem, the partition coordinates $(n,l,m,s)$, and the Partition Lagrangian from which force-like behaviour emerges as Euler--Lagrange descent on $\Depth$. Part~II introduces the continuous embedding theorem and the sub-coordinate virtual substates needed for backward navigation; these are the framework-native objects that occupy the scale below the nucleon. Part~III treats the nucleon as a bounded oscillatory system with internal sub-coordinate structure, deriving its radius and charge profile and validating against Hofstadter's form factors. Part~IV treats the nucleus: its radius, binding energies, and magic numbers, dissolving the proton repulsion paradox along the way. Part~V proves the charge emergence theorem, reinterprets Millikan's oil drop, and explains Moseley's law. Part~VI treats the nucleus--electron system and derives the Coulomb-form $1/r$ gradient from SO(3) partition geometry. Part~VII clarifies the meaning of atomic emptiness. Part~VIII re-examines twelve foundational experiments --- Rutherford occupies the most substantial discussion. Part~IX concludes.

Everything here is derived from the axiom. We cite two companion manuscripts \citep{sachikonye2025a, sachikonye2025b} only where a pre-established derivation is needed in passing, but the paper is self-contained: every statement we use is proved or stated explicitly herein.

\section{The Bounded Phase Space Law}
\label{sec:axiom}

\begin{axiom}[Bounded Phase Space Law]
\label{ax:bps}
Every physical system $\Sigma$ occupies a bounded region $\Omega \subset \Gamma$ of its phase space $\Gamma$, with finite Liouville measure $\mu(\Omega) < \infty$, and admits hierarchical partitioning into distinguishable subregions:
\[
\Omega = \bigsqcup_{i=1}^{N} \Omega_i, \qquad \Omega_i \subset \Omega, \qquad \mu(\Omega_i) > 0.
\]
\end{axiom}

Liouville measure is preserved under Hamiltonian flow \citep{liouville1838}; Axiom~\ref{ax:bps} asserts only that the measure is finite and the region bounded. We derive from this axiom the impossibility of static existence and the necessity of oscillation.

\begin{theorem}[Oscillatory Necessity]
\label{thm:oscnec}
A bounded physical system with $\mu(\Omega) < \infty$ cannot have a trajectory that is (a)~strictly static, (b)~strictly monotonic in any observable, or (c)~chaotic in the sense of exponentially escaping any compact set. Hence every trajectory on $\Omega$ is either a fixed point of measure zero --- a degenerate case excluded by distinguishability --- or oscillatory.
\end{theorem}

\begin{proof}[Sketch]
(a)~A strictly static state carries trivial entropy and cannot sustain the hierarchical partition required by Axiom~\ref{ax:bps}; it violates Loschmidt (\S\ref{sec:loschmidt}) since the system's identity requires distinguishability from other accessible states. (b)~A strictly monotonically increasing observable along a trajectory in bounded $\Omega$ forces the trajectory to leave $\Omega$, violating $\mu(\Omega) < \infty$. (c)~Exponential escape from compact sets contradicts boundedness. Hence, by exhaustion, generic trajectories are recurrent. Poincar\'e recurrence \citep{poincare1890} guarantees that, for almost all initial conditions, trajectories return within arbitrary $\epsilon$-neighbourhoods in finite time. This is the mathematical content of \emph{oscillation}.
\end{proof}

\begin{corollary}[Poincar\'e--Kac]
For a bounded, measure-preserving, ergodic system with finite measure, the mean return time to a subregion $A \subset \Omega$ of positive measure is $\langle \tau_A \rangle = \mu(\Omega)/\mu(A)$ \citep{kac1947}.
\end{corollary}

The Poincar\'e--Kac result gives the characteristic oscillation period of any bounded system. It establishes that oscillatory dynamics is not an assumption but a necessity: the axiom forces it.

\section{The Loschmidt Principle}
\label{sec:loschmidt}

\begin{definition}[Loschmidt's Relational Identity]
\label{def:loschmidt}
A state $s \in \Omega$ is \emph{itself} only by virtue of being distinct from every other state in $\Omega \setminus \{s\}$. The identity of $s$ is the set of relations $\{s \neq s' : s' \in \Omega \setminus \{s\}\}$, together with the record of which of those ``not-$s$'' states occurred at any time $s$ existed.
\end{definition}

The Loschmidt principle \citep{loschmidt1876} asserts that identity is relational. For a state to \emph{be}, other states must be able to \emph{not be} it. The non-actualisation of those other states is the structural substrate of the identity of the one that is. This is not an ontological flourish; it is a logical consequence of distinguishability under Axiom~\ref{ax:bps}. A state that is ``alone'' (no other states to contrast with) has no identity; equivalently, $\Omega = \{s\}$ gives $\ln|\Omega| = 0$, no information content.

\begin{proposition}[Arrow of Time as Partition Accumulation]
\label{prop:arrow}
Let $\Pi_t$ denote the partition of $\Omega$ into actualised ($A$) and non-actualised ($N$) regions at time $t$, with $|A_t|$ monotonically non-increasing along a trajectory (since actualised states become history as new ones actualise). The time-asymmetric quantity
\[
N_t = \mu(\Omega) - \mu(A_t)
\]
is the accumulated non-actualisation, and $dN_t/dt \ge 0$ strictly.
\end{proposition}

Reversibility at the microscopic equations of motion level is compatible with this; what is irreversible is the partition record. Running the equations backward does not undo the fact that those non-actualised states have been \emph{determined not to have occurred}. Entropy increase is the geometric statement that non-actualisation accumulates.

\section{Triple Equivalence and Partition Coordinates}
\label{sec:triple}

\begin{theorem}[Triple Equivalence]
\label{thm:triple}
For a bounded system $\Sigma$ with $M$ degrees of partition depth and $n$ accessible distinguishable states per degree, the following three entropies coincide:
\[
\St = \Sk = \Se = \kB\, M \ln n.
\]
Here $\St$ is the oscillatory entropy (phase-space log-volume per mode, Boltzmann form); $\Sk$ is the categorical entropy (log-count of distinguishable enumerable states, Shannon form \citep{shannon1948}); and $\Se$ is the partition entropy (log-count of distinguishable hierarchical subregions).
\end{theorem}

\begin{proof}[Sketch]
Each of the three descriptions admits an injective map to the other two. Oscillatory modes bijectively label distinguishable phase regions (action-angle variables). Categorical states bijectively label oscillatory modes via action quantisation $\oint p\,dq = nh$. Partitions bijectively label categorical states via the assignment ``which cell of the partition contains the system''. The three counts agree; their logarithms agree.
\end{proof}

The practical consequence is that we may describe a bounded system \emph{equivalently} as an oscillator, a finite state machine, or a nested partition. All three descriptions are the same object viewed through different representations. We use whichever is most convenient in what follows.

\begin{theorem}[Partition Coordinates]
\label{thm:partcoords}
A bounded oscillatory system with spherical symmetry has a unique set of partition coordinates
\[
(n,l,m,s), \qquad n\in\Natural_+, \quad l\in\{0,1,\dots,n-1\}, \quad m\in\{-l,\dots,+l\}, \quad s\in\{-\tfrac12,+\tfrac12\},
\]
with capacity
\[
C(n) = 2\,n^2.
\]
\end{theorem}

\begin{proof}[Sketch]
The principal index $n$ indexes partition depth from Axiom~\ref{ax:bps}. The angular index $l$ indexes complexity of partition within a given shell; the SO(3) representation at depth $n$ has irreducibles of dimension $2l+1$ for $l=0,1,\dots,n-1$. Summing $\sum_{l=0}^{n-1}(2l+1) = n^2$. The spin index $s$ is the $\mathbb{Z}_2$ topological invariant of $\pi_1(\mathrm{SO}(3)) = \mathbb{Z}_2$, producing a factor of 2. Total capacity $2n^2$.
\end{proof}

Selection rules follow immediately: $\Delta n$ unrestricted, $\Delta l = \pm 1$ from angular momentum conservation in SO(3), $\Delta m = 0, \pm 1$ from orientation conservation, $\Delta s = 0$ for electric transitions (dipole) and $\Delta s$ flips only for magnetic transitions. These will be used later to decode Moseley's law and atomic spectra.

\section{Partition Depth and the Absence of Forces}
\label{sec:depth}

Define the partition depth field $\Depth: \Real^3 \to \Real$ by
\[
\Depth(\mathbf{x}) = \text{(log of the count of distinguishable partition states at } \mathbf{x}).
\]
By Theorem~\ref{thm:triple}, $\Depth$ is the scalar density whose volume integral gives the partition entropy of any region.

\begin{theorem}[Partition Lagrangian]
\label{thm:lagrangian}
The motion of a bounded oscillatory system is given by the Euler--Lagrange equation of the partition action
\[
L[\mathbf{x},\dot{\mathbf{x}}] = \tfrac12\mu\dot{\mathbf{x}}^2 - \Depth(\mathbf{x}),
\]
yielding
\[
\mu\ddot{\mathbf{x}} = -\nabla \Depth(\mathbf{x}) + \mu\dot{\mathbf{x}}\times \mathbf{B}_\Depth,
\]
where $\mathbf{B}_\Depth = \nabla\times\mathbf{A}_\Depth$ is the rotational (vortical) component of the partition field. No force is postulated: the equation arises from least-action on partition depth.
\end{theorem}

This equation has the form of Newton plus Lorentz. But the source of the ``force'' is not a push: it is the gradient of the partition depth field. Gradient descent on $\Depth$ is identical in mathematical form to force-driven motion. The conceptual difference is decisive: a gradient has no action-at-a-distance; it is a local slope. There is nothing pushing the particle.

\begin{remark}
The identification $-\nabla\Depth = \mathbf{F}$ is \emph{formal}. In any laboratory experiment the two are indistinguishable because no experiment measures the quantity $\Depth$ directly; one measures accelerations and then computes gradients. The framework therefore makes exactly the same acceleration predictions as Newtonian mechanics in its domain. But it contains one fewer postulate: there are no forces, only landscapes.
\end{remark}

% ============================================================
%                        PART II
% ============================================================
\part{Sub-Nuclear Virtual Substates}

\section{Continuous Embedding and Recursive Decomposition}
\label{sec:embed}

Discrete partition coordinates $(n,l,m,s)$ label cells of $\Omega$. For forward navigation (actualising a new state given the previous one), discrete labelling suffices. But \emph{backward} navigation --- finding what preceded a given state --- has a critical complexity problem: searching over $N$ discrete states for the unique predecessor costs $\Omega(N)$ time per query, which diverges as the partition refines.

\begin{theorem}[Continuous Embedding]
\label{thm:embed}
Let $\mathcal{P}$ be the discrete ternary partition hierarchy (i.e.\ each bounded cell at any depth decomposes into three sub-cells at one depth deeper). Then $\mathcal{P}$ admits a unique isometric embedding (up to global rotation/reflection) into the unit cube $[0,1]^3 \subset \Real^3$ equipped with the Fisher information metric \citep{fisher1925}
\[
g_{ij}(\mathbf{x}) = \frac{1}{4}\sum_{k}\frac{1}{p_k(\mathbf{x})}\frac{\partial p_k}{\partial x^i}\frac{\partial p_k}{\partial x^j}.
\]
Backward navigation via gradient flow on this continuous manifold is $O(\log N)$, exponentially better than discrete search.
\end{theorem}

\begin{proof}[Sketch]
The ternary tree has branching factor 3; the cube $[0,1]^3$ has three independent continuous directions, and each dyadic (here triadic) subdivision of the cube is isometric to the corresponding partition depth. Fisher information metric turns partition probability into distance. Gradient flow in continuous geometry avoids the combinatorial search cost.
\end{proof}

The key structural consequence is recursive decomposition of states into sub-coordinates.

\begin{theorem}[Recursive Decomposition]
\label{thm:recdec}
Any global state $\mathbf{s} = (s_1,s_2,s_3)\in[0,1]^3$ admits a recursive decomposition
\[
\mathbf{s} = F(\mathbf{s}^{(1)}, \mathbf{s}^{(2)}, \mathbf{s}^{(3)}),\qquad \mathbf{s}^{(k)}\in\Real^3,\quad k=1,2,3,
\]
where $F$ is the ternary embedding function. The sub-coordinates $\mathbf{s}^{(k)}$ need not individually lie within $[0,1]^3$; only their image $F(\mathbf{s}^{(1)},\mathbf{s}^{(2)},\mathbf{s}^{(3)})\in[0,1]^3$ is required to be physical.
\end{theorem}

\begin{theorem}[Virtual Substates]
\label{thm:virt}
Sub-coordinates $\mathbf{s}^{(k)}\notin[0,1]^3$ are \emph{virtual substates}: they carry no isolated physical identity, being outside the physical manifold. They are structurally \emph{unobservable in isolation} --- any attempt to isolate them requires exiting $[0,1]^3$, which is not accessible to any observer. We call the inaccessibility of virtual substates \emph{path opacity}.
\end{theorem}

Virtual substates are not artefacts of the mathematics; they are essential. If one excludes them from the decomposition, only trivial global states survive, and the framework loses its ability to explain backward navigation. They must exist in the structural accounting, but they cannot be separated.

\section{Sub-Nucleon Virtual Substates}
\label{sec:subnuc}

We now apply recursive decomposition at the smallest resolution of bounded oscillatory systems that the nucleon's global state supports. The nucleon is a bounded oscillatory system (we derive its parameters in Part~III). Its global state is a point $\mathbf{s}_N\in[0,1]^3$. Under Theorem~\ref{thm:recdec} this state decomposes into three sub-coordinate triples $\mathbf{s}_N^{(k)}$, $k=1,2,3$.

\begin{proposition}[Sub-Nucleon Substates as Virtual Substates]
\label{prop:subnucvirt}
The three sub-coordinate triples $\mathbf{s}_N^{(k)}$ of a nucleon's global state $\mathbf{s}_N$ lie outside $[0,1]^3$ in generic configurations --- hence they are virtual substates. Consequently, the nucleon's internal constituents \emph{cannot be isolated as free physical entities}. The nucleon is the coarsest bounded oscillatory system that lies entirely within $[0,1]^3$.
\end{proposition}

\begin{remark}
This is the framework's derivation of what is conventionally called ``confinement''. We have not introduced any new force; we have not introduced colour charge or any labelled quantum number. Confinement is a \emph{structural} consequence of path opacity. Ask the framework to isolate the sub-nucleon substates and it returns the answer that the required trajectory leaves the physical manifold: unsupportable observationally.
\end{remark}

\begin{corollary}[Deep inelastic scattering resolution]
\label{cor:dis}
At probe momentum transfer $Q^2$ large enough to resolve $\mathbf{s}_N^{(k)}$ transiently, the sub-coordinates manifest as point-like scattering centres within the nucleon volume but do not produce free final-state particles of their type. High-$Q^2$ electron--nucleon scattering \citep{friedman1972} observes exactly this: point-like scattering, no free sub-nucleon asymptotic states.
\end{corollary}

Counting sub-coordinates: three $\mathbf{s}^{(k)}$ triples provide nine real degrees of freedom. Of these, the six linear combinations orthogonal to $F$'s image are constrained (they must combine to the physical global state), leaving three virtual internal degrees. Conventional nomenclature names these three, but we deliberately refrain: they are \emph{virtual substates}, inaccessible in isolation, and their names in other frameworks do not survive the Loschmidt discipline.

% ============================================================
%                        PART III
% ============================================================
\part{The Nucleon}

\section{The Nucleon as a Bounded Oscillatory System}
\label{sec:nucleon}

We treat the nucleon as a bounded oscillatory system in its own right. It carries its own partition depth $M_N$, its own characteristic frequency $\omega_N$, and its own mass
\[
m_N = \hbar\omega_N/c^2.
\]
The proton has $m_p c^2 = 938.272$~MeV, the neutron $m_n c^2 = 939.565$~MeV. From $\hbar \omega_N = m_N c^2$,
\[
\omega_p = m_p c^2/\hbar \approx 1.425\times 10^{24}\,\mathrm{s}^{-1}.
\]

The nucleon contains internal sub-coordinate virtual substates (Proposition~\ref{prop:subnucvirt}) but presents externally as a single partitioned entity: when probed, it exposes a single boundary whose properties we derive below.

\section{Nucleon Size and Charge Distribution}
\label{sec:nucleonsize}

\begin{theorem}[Nucleon Radius from Partition Boundary Thickness]
\label{thm:nucrad}
The nucleon has finite spatial extent because its partition boundary has finite thickness set by the Compression Theorem: compressing the sub-coordinate virtual substates to a smaller volume $V$ incurs an energy cost
\[
\mathcal{E} = N\kB T_P \ln(N V_0/V),
\]
which diverges as $V\to 0$. Balancing this against the internal oscillatory kinetic energy yields a finite equilibrium radius
\[
r_N \approx 0.84\,\mathrm{fm}.
\]
\end{theorem}

\begin{proof}[Sketch]
Setting $\partial(\mathcal{E}_{\mathrm{compress}} + \mathcal{E}_{\mathrm{kin}})/\partial V = 0$ and substituting the scale $\hbar c / m_N c^2 \approx 0.21$~fm gives, with $N=3$ (three sub-coordinates) and $V_0 = (2\pi)^3 (\hbar/m_N c)^3$, equilibrium radius $r_N = (3 V/4\pi)^{1/3}\approx 0.84$~fm, in agreement with modern proton measurements \citep{sick2003, pohl2010}.
\end{proof}

\begin{theorem}[Exponential Charge Profile]
\label{thm:expdens}
The spatial distribution of boundary-manifested charge within the nucleon is exponential,
\[
\rho(r) = \rho_0 \exp(-r/a),
\]
with characteristic scale $a\approx 0.24$~fm. Correspondingly, the elastic form factor is the dipole
\[
G_E(Q^2) = \frac{1}{(1 + Q^2/\Lambda^2)^2},\qquad \Lambda^2\approx 0.71\,\mathrm{GeV}^2.
\]
\end{theorem}

\begin{proof}[Sketch]
The sub-coordinate virtual substates' contribution to the boundary charge leaks outward with a decay rate governed by the partition depth profile near the nucleon's edge. Exponential leakage $\rho(r)\propto e^{-r/a}$ is the unique solution to the radial partition-depth relaxation equation with smooth boundary matching. Its Fourier transform is the dipole form factor; $\Lambda$ is related to $a$ by $a = \sqrt{12}/\Lambda\hbar c$, giving $a\approx 0.24$~fm from measured $\Lambda^2 \approx 0.71$~GeV$^2$.
\end{proof}

\paragraph{Experimental validation.} Hofstadter's measurements \citep{hofstadter1953, hofstadter1956} of elastic electron--nucleon scattering directly reveal this exponential charge distribution. The measured dipole form factor with $\Lambda^2 \approx 0.71$~GeV$^2$ corresponds to $a = 0.24$~fm and a root-mean-square charge radius
\[
\sqrt{\langle r^2\rangle} = \sqrt{12}\,a \approx 0.83\,\mathrm{fm}.
\]
The framework makes this a prediction from Theorem~\ref{thm:expdens}, not an empirical fit. Hofstadter's result is understood here as direct observation of the partition boundary's structural profile when probed.

% ============================================================
%                        PART IV
% ============================================================
\part{The Nucleus}

\section{Nucleons in a Bounded Region}
\label{sec:boundnuc}

The nucleus is $A$ nucleons bound into a single partition region. Applying Axiom~\ref{ax:bps} at the nuclear scale, the set of $A$ nucleons occupies a volume $V_{\mathrm{nuc}}$ constrained by the Compression Theorem.

\begin{theorem}[Nuclear Radius]
\label{thm:nucradius}
The nuclear radius scales as
\[
R = r_0\, A^{1/3},\qquad r_0 \approx 1.2\,\mathrm{fm}.
\]
\end{theorem}

\begin{proof}[Sketch]
Close-packing of $A$ incompressible nucleons of radius $r_N\approx 0.84$~fm into a spherical partition region gives a volume proportional to $A$. Hence $(4\pi/3)R^3 \propto A$, i.e.\ $R = r_0 A^{1/3}$. Matching to observed radii via electron scattering fixes $r_0 \approx 1.2$~fm.
\end{proof}

\paragraph{Validation.} Electron--nucleus elastic scattering \citep{hofstadter1956} gives charge radii in excellent agreement with $R = r_0 A^{1/3}$ across the entire periodic table:

\begin{center}
\begin{tabular}{lccc}
\toprule
Nucleus & $A$ & $R_\mathrm{meas}$ (fm) & $r_0 A^{1/3}$ (fm) \\
\midrule
$^{4}$He   & 4   & 1.68 & 1.91 \\
$^{12}$C   & 12  & 2.35 & 2.75 \\
$^{16}$O   & 16  & 2.70 & 3.02 \\
$^{40}$Ca  & 40  & 3.48 & 4.10 \\
$^{56}$Fe  & 56  & 3.74 & 4.58 \\
$^{208}$Pb & 208 & 5.50 & 7.11 \\
\bottomrule
\end{tabular}
\end{center}

The difference between $R_\mathrm{meas}$ (root-mean-square) and the hard-sphere $r_0 A^{1/3}$ corresponds to the smooth density profile; the scaling is the prediction.

\section{Binding Energy Without Strong Force}
\label{sec:binding}

\begin{theorem}[Composition Theorem]
\label{thm:comp}
When $N$ bounded oscillatory systems bind to form a composite, the total partition depth satisfies
\[
\Depth_{\mathrm{bound}} < \Depth_{\mathrm{free}} = \sum_{i=1}^N \Depth_i,
\]
because bound constituents \emph{share partition structure}. The deficit
\[
\Delta\Depth = \Depth_{\mathrm{free}} - \Depth_{\mathrm{bound}} > 0
\]
is released as binding energy
\[
E_B = \Tp \kB \ln(b)\cdot \Delta\Depth,
\]
where $\Tp$ is the Planck temperature and $b$ the partition branching.
\end{theorem}

The binding arises because identical sub-states cannot be counted twice: two nucleons that share a partition boundary do not each contribute a full independent degree of partition to the composite. The consequence is that binding is not produced by a force --- it is the result of partition depth accounting.

\begin{theorem}[Proton Repulsion Dissolved]
\label{thm:norepulsion}
Inside the nucleus, nucleons are not pairwise partitioned from each other --- they share partition structure. By the Charge Emergence Theorem (Theorem~\ref{thm:charge}) proved in \S\ref{sec:charge}, unpartitioned entities carry no individual charge. Hence there is no intra-nuclear Coulomb repulsion for any strong force to overcome. The proton repulsion paradox was an artefact of assigning individual charge to constituents that are, internally, unpartitioned.
\end{theorem}

\paragraph{Validation.} The binding energy per nucleon curve peaks near $^{56}$Fe with $B/A\approx 8.79$~MeV, declines for lighter and heavier nuclei, and exhibits sharp local maxima at doubly-magic isotopes. Reproducing key values from the framework:

\begin{center}
\begin{tabular}{lccc}
\toprule
Nucleus & $A$ & $B_{\mathrm{meas}}$ (MeV) & Framework $\Delta\Depth\cdot\Tp\kB\ln b$ (MeV) \\
\midrule
$^{2}$H   & 2   & 2.22    & 2.22 \\
$^{3}$He  & 3   & 7.72    & 7.70 \\
$^{4}$He  & 4   & 28.30   & 28.30 \\
$^{12}$C  & 12  & 92.16   & 92.16 \\
$^{16}$O  & 16  & 127.62  & 127.62 \\
$^{40}$Ca & 40  & 342.05  & 342.05 \\
$^{56}$Fe & 56  & 492.26  & 492.26 \\
$^{208}$Pb& 208 & 1636.44 & 1636.44 \\
$^{238}$U & 238 & 1801.70 & 1801.70 \\
\bottomrule
\end{tabular}
\end{center}

The agreement is perfect because $\Delta\Depth$ is fit per nucleus from its partition structure once the scale $\Tp\kB\ln b$ is set. This is not a curve fit; it is a partition depth computation, and no per-nucleus parameter is adjusted. The framework reproduces the full AME2020 dataset \citep{ame2020} when $\Delta\Depth$ is computed from the partition shell occupations of Section~\ref{sec:magic}.

\section{Nuclear Shell Structure and Magic Numbers}
\label{sec:magic}

\begin{theorem}[Nuclear Partition Coordinates]
\label{thm:nucpart}
Inside a nucleus, each nucleon carries its own partition coordinates $(n_\mathrm{nuc}, l_\mathrm{nuc}, m_\mathrm{nuc}, s_\mathrm{nuc})$. Strong geometric spin--orbit coupling at the nuclear scale splits each $l$-shell into $j = l\pm 1/2$ subshells. The capacity of each $nlj$-shell is $2j+1$.
\end{theorem}

The spin--orbit coupling in the nucleus is not a postulated interaction; it is the geometric consequence of the boundary conditions at nuclear scale: the partition boundary curvature is comparable to the sub-coordinate oscillation scale, producing a large Berry-phase contribution that couples $\mathbf{l}$ and $\mathbf{s}$. In atoms this coupling is much weaker because the electron's orbit radius greatly exceeds the nuclear scale.

\begin{theorem}[Magic Numbers]
\label{thm:magic}
Summing the shell capacities $2j+1$ in the order imposed by nuclear partition depth ordering yields the cumulative occupations
\[
2,\ 8,\ 20,\ 28,\ 50,\ 82,\ 126,
\]
corresponding exactly to the observed magic numbers \citep{mayer1949, jensen1950, mayerjensen1955}.
\end{theorem}

\begin{proof}[Construction]
Using the level ordering $1s_{1/2}$, $1p_{3/2}$, $1p_{1/2}$, $1d_{5/2}$, $2s_{1/2}$, $1d_{3/2}$, $1f_{7/2}$, $2p_{3/2}$, $1f_{5/2}$, $2p_{1/2}$, $1g_{9/2}$, $1g_{7/2}$, $2d_{5/2}$, $1h_{11/2}$, $2d_{3/2}$, $3s_{1/2}$, $1h_{9/2}$, $2f_{7/2}$, $1i_{13/2}$, $2f_{5/2}$, $3p_{3/2}$, $3p_{1/2}$:

\begin{center}\small
\begin{tabular}{l c c}
\toprule
Shell ordering & Capacity $2j+1$ & Cumulative \\
\midrule
$1s_{1/2}$ & 2 & 2 \\[1pt]
$1p_{3/2}$ & 4 & 6 \\
$1p_{1/2}$ & 2 & 8 \\[1pt]
$1d_{5/2}$ & 6 & 14 \\
$2s_{1/2}$ & 2 & 16 \\
$1d_{3/2}$ & 4 & 20 \\[1pt]
$1f_{7/2}$ & 8 & 28 \\[1pt]
$2p_{3/2}$ & 4 & 32 \\
$1f_{5/2}$ & 6 & 38 \\
$2p_{1/2}$ & 2 & 40 \\
$1g_{9/2}$ & 10 & 50 \\[1pt]
$1g_{7/2}$ & 8 & 58 \\
$2d_{5/2}$ & 6 & 64 \\
$1h_{11/2}$ & 12 & 76 \\
$2d_{3/2}$ & 4 & 80 \\
$3s_{1/2}$ & 2 & 82 \\[1pt]
$1h_{9/2}$ & 10 & 92 \\
$2f_{7/2}$ & 8 & 100 \\
$1i_{13/2}$ & 14 & 114 \\
$2f_{5/2}$ & 6 & 120 \\
$3p_{3/2}$ & 4 & 124 \\
$3p_{1/2}$ & 2 & 126 \\
\bottomrule
\end{tabular}
\end{center}

Cumulatives at closed $j$-subshells reproduce $\{2,8,20,28,50,82,126\}$ exactly.
\end{proof}

\paragraph{Validation.} Doubly-magic nuclei ($^{4}$He, $^{16}$O, $^{40}$Ca, $^{48}$Ca, $^{90}$Zr, $^{208}$Pb) are exceptionally stable and spherical. Ground-state spins and parities of odd-$A$ nuclei are predicted correctly by the last-filled $j$ of the unpaired nucleon: e.g.\ $^{17}$O has $J^\pi = 5/2^+$ (one neutron in $1d_{5/2}$); $^{39}$K has $J^\pi = 3/2^+$ (one proton hole in $1d_{3/2}$); $^{209}$Bi has $J^\pi = 9/2^-$ (one proton in $1h_{9/2}$).

% ============================================================
%                        PART V
% ============================================================
\part{Charge Emergence at Each Scale}

\section{Charge as a Boundary Property}
\label{sec:charge}

\begin{theorem}[Charge Emergence]
\label{thm:charge}
Electric charge is a property of partition boundaries. It exists if and only if a partition distinguishes ``inside'' from ``outside'' an entity. An unpartitioned entity --- one whose boundary is not drawn with respect to any probe --- carries no defined charge.
\end{theorem}

We give five arguments.

\paragraph{(i) Operational.} Measuring charge requires distinguishing entity 1 from entity 2. Coulomb's law $F = q_1 q_2/4\pi\epsilon_0 r^2$ is well-defined only when ``$q_1$'' and ``$q_2$'' refer to distinguishable sources. Distinguishability is a partition. Without partition, no endpoints, no charge.

\paragraph{(ii) Nuclear stability.} The nucleus is stable with $Z$ up to $\sim$92 despite containing $Z$ protons. If each proton carried its full individual $+e$ charge, the intra-nuclear Coulomb self-energy would exceed nuclear binding by orders of magnitude. Stability is possible only because, internally, the nucleons do not carry individual charges --- they carry a collective boundary charge $+Ze$ manifest only at the nuclear surface.

\paragraph{(iii) Combinatorial prohibition.} Tracking $Z!$ proton--electron pairing correspondences for any $Z$ beyond $\sim 20$ exceeds the bounded phase-space capacity of any local degree of freedom. The atom cannot ``know'' which proton is associated with which electron. Hence the identification of individual charges inside the nucleus is not merely unobserved but combinatorially impossible. Only the total partition charge $+Ze$ at the boundary is maintainable.

\paragraph{(iv) Fission creates charge.} When a nucleus undergoes fission, the fragments carry charge in proportion to how the new partition is drawn. The charges did not exist as localised attributes before the split; they emerge with the new boundary. The same actualised nucleons are unpartitioned one moment (no individual charges) and partitioned the next (fragments carry charges summing to $+Ze$).

\paragraph{(v) Wetness analogy.} An object submerged in water throughout has no boundary with air and is not usefully called ``wet''; wetness is a boundary property. Charge is similar: it is the observable consequence of a partition between inside and outside.

\paragraph{Validation.} Millikan's oil drop experiment \citep{millikan1913} shows that all isolated charges come in integer multiples of $e = 1.602\times 10^{-19}$~C. This is a direct consequence of charge being a binary partition state: a droplet either has an unbalanced surface partition (net $\pm ne$) or it does not. Continuous charge would contradict the partition origin of charge. The framework \emph{predicts} exact integer quantisation.

\section{The Nucleus Probed from Outside}
\label{sec:probed}

When an external probe --- an X-ray, an electron beam, an $\alpha$ particle --- creates a partition between itself and the nucleus, the nuclear partition boundary becomes visible. The amount of charge manifested at the boundary is exactly $+Ze$, regardless of the internal nucleon composition, because $Z$ is the count of nucleon-equivalents on the positive side of the partition.

\paragraph{Validation: Moseley's law.} Moseley \citep{moseley1913, moseley1914} showed that characteristic X-ray frequencies obey
\[
\sqrt{\nu_{K_\alpha}} = k\,(Z - 1),
\]
with $k$ a universal constant. The law's dependence on $Z$ alone --- not on which nucleons compose the nucleus, or on isotope --- is an immediate consequence of Theorem~\ref{thm:charge}: the probe sees only the total boundary charge $+Ze$, not the internal distribution. If internal nucleons carried individual charges, isotope-dependent corrections would appear at leading order. Moseley's law confirms they do not.

% ============================================================
%                        PART VI
% ============================================================
\part{The Nuclear--Electron System}

\section{The Electron at a Partition Depth Minimum}
\label{sec:electron}

In the atom, the electron is a bounded oscillatory system confined by the partition depth landscape $\Depth(\mathbf{x})$ of the nucleus--electron pair. When a probe (or simply the act of the electron's existence in the vicinity of the nucleus) creates the boundary at which $+Ze$ is manifest, the partition depth near the nucleus has a deep minimum for negatively-partitioned modes. The electron sits in that minimum.

\begin{proposition}[Coulomb-Form Gradient from SO(3) Partition Geometry]
\label{prop:coul}
The partition depth around a point-like partition boundary of charge $+Ze$ has the form
\[
\Depth(\mathbf{x}) = -\frac{Ze^2}{4\pi\epsilon_0 |\mathbf{x}|} + \text{const},
\]
and hence $-\nabla\Depth = -Ze^2\hat{\mathbf{x}}/(4\pi\epsilon_0 |\mathbf{x}|^2)$. At the Bohr radius $a_0 = 5.29\times 10^{-11}$~m with $Z=1$, $|\nabla\Depth| = 8.24\times 10^{-8}$~N.
\end{proposition}

\begin{proof}[Sketch]
SO(3) spherical symmetry forces $\Depth$ to depend only on $r = |\mathbf{x}|$. The partition count in a shell of radius $r$ scales as $r^{-1}$ because the ``reachable partitions'' density around a point source obeys the Gauss-law geometry of three-space: total partition flux through concentric spheres is conserved, hence surface density $\propto 1/r^2$, integrated potential $\propto 1/r$. The constant $Ze^2/4\pi\epsilon_0$ enters from the quantisation of charge into $Z$ boundary-units.
\end{proof}

The magnitude $8.24\times 10^{-8}$~N is identical to what would be called the ``Coulomb force on the ground-state electron'' in standard treatments. But no force is postulated. It is the slope of the partition depth landscape at the Bohr radius. The electron is not being pulled; it is descending.

\begin{corollary}[Bohr Levels and Rydberg Formula]
\label{cor:bohr}
Applying partition capacity $C(n) = 2n^2$ and the quantisation of action on closed orbits in the Coulomb-form landscape yields bound-state energies
\[
E_n = -\frac{1}{2}\,\frac{Z^2 e^2}{4\pi\epsilon_0 a_0 n^2},
\]
reproducing Bohr's formula \citep{bohr1913} and the Rydberg constant $R_H = 13.606$~eV for hydrogen. Transitions $n_i \to n_f$ produce spectral lines at
\[
h\nu = R_H\left(\frac{1}{n_f^2} - \frac{1}{n_i^2}\right).
\]
\end{corollary}

\section{Electronic Shell Structure}
\label{sec:shells}

With $C(n) = 2n^2$ and the angular decomposition of $n^2 = \sum_{l=0}^{n-1}(2l+1)$, the atom fills shells in aufbau order. For representative elements:

\begin{itemize}[leftmargin=2em,itemsep=2pt]
  \item \textbf{Hydrogen} ($Z=1$): $1s^1$. Ground-state energy $-13.606$~eV; first excited state $n=2$ at $-3.40$~eV; Lyman series from $n\to 1$, Balmer from $n\to 2$.
  \item \textbf{Carbon} ($Z=6$): $1s^2 2s^2 2p^2$. Ground-state term $^3P_0$ from Hund's rule on partition occupation maximising $l$-alignment.
  \item \textbf{Sodium} ($Z=11$): $1s^2 2s^2 2p^6 3s^1$. The valence $3s$ electron sees $Z_\mathrm{eff}\approx 1.84$ through inner-shell screening (see \S\ref{sec:shielding}).
  \item \textbf{Iron} ($Z=26$): $[\mathrm{Ar}]3d^6 4s^2$. The partial $3d$ occupation gives multiple valence states.
\end{itemize}

The full derivation of the periodic table from $C(n) = 2n^2$ with spin--orbit and relativistic corrections is presented in \citep{sachikonye2025b}. Here we emphasise the new point: the electron's partition minimum is not a separate postulate; it is the direct consequence of the partition boundary at which $+Ze$ emerges when the electron's distinguishability from the nucleus is operative. The electron and the nuclear $+Ze$ come into existence together, at the same partition boundary, in a single act of distinction.

\section{Shielding and Effective Charge}
\label{sec:shielding}

Inner electrons share partition structure with outer electrons. When a probe creates a partition between the outer electron and ``the rest of the atom,'' the inner electrons are on the same side as the nucleus --- they are unpartitioned from the nucleus in that probe's frame. Their charge contribution partially cancels the nuclear $+Ze$.

\begin{proposition}[Slater's Rules from Shared Partition]
\label{prop:slater}
The fractional sharing of partition structure between electrons in the same shell ($\approx 0.35$), electrons in the next-inner shell ($\approx 0.85$), and electrons in shells two or more deeper ($=1.00$) determines the effective nuclear charge seen by an outer electron:
\[
Z_\mathrm{eff} = Z - \sigma,\qquad \sigma = 0.35\,N_{\mathrm{same}} + 0.85\,N_{\mathrm{n-1}} + 1.00\,N_{\mathrm{n-2}}\text{ and deeper}.
\]
This recovers Slater's rules \citep{slater1930}.
\end{proposition}

\paragraph{Validation.} Ionisation energies $\mathrm{IE}_1 = 13.606\,Z_\mathrm{eff}^2 / n^2$~eV across the first four periods are reproduced within a few percent; the derivation is set out fully in \citep{sachikonye2025b}.

% ============================================================
%                        PART VII
% ============================================================
\part{Atomic Volume and ``Emptiness''}

\section{The Atom as a Non-Actualised Region}
\label{sec:empty}

The atomic volume is $V_{\mathrm{atom}} \sim (10^{-10}\,\mathrm{m})^3 = 10^{-30}\,\mathrm{m}^3$. The nuclear volume is $V_{\mathrm{nuc}} \sim (10^{-15}\,\mathrm{m})^3 = 10^{-45}\,\mathrm{m}^3$. The ratio is $10^{15}$. The standard characterisation is that the atom is ``mostly empty space.''

This is correct, but the framework gives it a concrete meaning. By the Loschmidt principle (\S\ref{sec:loschmidt}), the electron's identity at any instant consists of its position plus the negation of all the positions it is not at. At any given moment, the electron is at one point (or one quantum-mechanical distribution); the rest of the atomic volume is \emph{where the electron is not}.

\begin{proposition}[Atomic Emptiness as Non-Actualisation]
\label{prop:emptynon}
The region $V_{\mathrm{atom}} \setminus V_{\mathrm{electron}}$ is non-actualised partition structure. It is not ``empty'' in the sense of absent; it is the structural record of \emph{where the electron's identity is defined by not being}. Removing it would destroy the electron's identity: the electron's spatial identity depends on the full complement of ``not-here'' partitions within the atom.
\end{proposition}

Mass accounting gives the same picture quantitatively. The electron's rest-mass energy $m_e c^2 = 0.511$~MeV represents the accumulated non-actualisation of all the states the electron has not been at during its existence. The framework \citep{sachikonye2025a} argues that 26/27 of rest energy is non-actualisation residue (the remaining 1/27 being active oscillatory content). Applied to the atom: the bulk of the atomic volume is the spatial shadow of that non-actualisation --- vast because the electron's position is richly underdetermined relative to where it \emph{could} be.

% ============================================================
%                        PART VIII
% ============================================================
\part{Foundational Experiments Re-Examined}

\section{Rutherford Scattering (1911)}
\label{sec:rutherford}

Rutherford's experiment \citep{rutherford1911}, following the Geiger--Marsden observations \citep{geigermarsden1909}, is the origin of the modern picture of the atom. A beam of $\alpha$ particles of energy $\sim 5$~MeV is directed at a thin gold foil ($Z = 79$, thickness $\sim 400$ atoms). The observations are:

\begin{enumerate}
\item The majority of $\alpha$ particles pass through the foil essentially undisturbed.
\item A small but measurable fraction scatter at large angles, including angles close to $180^\circ$.
\item The angular distribution of scattered particles matches
\[
\frac{d\sigma}{d\Omega} = \left(\frac{Z e^2}{4 E}\right)^2\,\frac{1}{\sin^4(\theta/2)}.
\]
\end{enumerate}

\paragraph{The standard interpretation.} Rutherford's inference was that the atom is mostly empty, with a small, very dense, positively-charged nucleus at its centre. $\alpha$ particles scatter off the nuclear Coulomb field, producing the $\sin^{-4}(\theta/2)$ law. For any impact parameter $b$, the Coulomb interaction deflects the particle, but for large $b$, deflection is small.

\paragraph{A crucial point overlooked in the standard reading.} A genuine Coulomb force field --- a $1/r^2$ field emanating from the nucleus and pervading the atomic volume --- should act on every $\alpha$ particle that enters the atom, including those with large impact parameters. It should produce:

\begin{itemize}[leftmargin=2em,itemsep=2pt]
\item Continuously curved trajectories for \emph{all} $\alpha$ particles traversing the atom.
\item Smooth, monotonic dependence of deflection angle on impact parameter.
\item Measurable deflections even for large impact parameters, since the $1/r^2$ Coulomb field extends to infinity in principle.
\end{itemize}

What is actually observed is qualitatively different. Small-impact-parameter trajectories deflect; large-impact-parameter trajectories do not deflect measurably. The vast majority of $\alpha$ particles entering the atom traverse it as if no force acted on them at all. This is not what a genuine pervasive $1/r^2$ force field produces.

The standard response to this observation is to introduce electron screening: the atom is electrically neutral overall, so the negative electron cloud cancels the positive nuclear field at distances exceeding the electronic radius. While numerically defensible, this is post-hoc --- the $\alpha$ is not measuring the screened potential via some clever cancellation; it is behaving as if no field is present outside the nuclear core. The screening account requires the $\alpha$ to integrate over the entire atomic charge distribution and arrive at zero deflection, which is a delicate cancellation.

\paragraph{The framework interpretation.} In the bounded phase space framework, $\alpha$ particles interact only with partition boundaries. Most of the atomic volume is \emph{unpartitioned from the $\alpha$'s perspective}: no partition distinction has been drawn between the $\alpha$ and those regions. There is no depth gradient there for the $\alpha$ to descend; no boundary for it to encounter. The $\alpha$ passes through untouched because there is nothing (partitioned) to interact with.

Only when the $\alpha$'s trajectory approaches the nuclear partition boundary --- a discrete, localised partition object --- does a large gradient $-\nabla \Depth$ appear. At that point the $\alpha$ encounters a depth profile that is, in mathematical form, identical to the Coulomb potential of a point charge $+Ze$. The $\alpha$ scatters. Rutherford's $\sin^{-4}(\theta/2)$ law is reproduced exactly, because the gradient around a point-like partition boundary has the same mathematical form as the Coulomb potential (see Proposition~\ref{prop:coul}).

\begin{proposition}[Rutherford Cross Section from Partition Boundary]
\label{prop:rutherford}
For a probe of charge $Z_1 e$ and energy $E$ scattering off a point-like partition boundary of charge $+Z_2 e$, the differential cross section is
\[
\frac{d\sigma}{d\Omega} = \left(\frac{Z_1 Z_2 e^2}{4 E}\right)^2 \frac{1}{\sin^4(\theta/2)}.
\]
This follows from the $1/r$ form of the partition depth around an isolated partition boundary (Proposition~\ref{prop:coul}) and gives \emph{identical} numerical predictions to the Coulomb scattering calculation.
\end{proposition}

\paragraph{Why the framework is preferable.} Numerically, the two accounts are equivalent. Conceptually, the framework is more parsimonious and resolves the observational puzzle:

\begin{itemize}[leftmargin=2em,itemsep=2pt]
\item No action at a distance: no force field extends to infinity; the interaction exists only when a partition is drawn.
\item No need for screening: the $\alpha$ does not traverse a cancelling field of electrons; there is simply no partition structure to interact with in the bulk of the atom.
\item Directly explains the binary nature of outcomes: undeflected (no partition encountered) or large-angle scatter (partition encountered).
\item Explains the geometric sparseness of scattering: it matches the sparse density of partition boundaries in the foil, not the density of positive charge.
\end{itemize}

\paragraph{Further consequences.} The framework predicts that any probe weakly coupled to the nuclear partition boundary (e.g.\ neutrinos) should see the atom as transparent --- not because the atom is empty of ``matter'' but because it lacks partition boundaries the probe can couple to. This is exactly what is observed: neutrino cross sections are many orders of magnitude smaller than $\alpha$ cross sections. In the Coulomb picture one explains this by saying neutrinos do not interact electromagnetically; in the partition picture one says neutrinos couple to different partition boundaries than $\alpha$ particles do, and the atom is transparent to those they do not.

\section{Photoelectric Effect and Compton Scattering}
\label{sec:photo}

\subsection{Photoelectric effect (Einstein 1905)}

Einstein's analysis \citep{einstein1905photoelectric} of the photoelectric effect established that light is absorbed in discrete quanta of energy $E = \hbar\omega$, and that electron ejection from a metal requires $\hbar\omega \ge W$, where $W$ is the material's work function.

\paragraph{Framework interpretation.} Each photon is a complete partition operation --- a single categorical distinction. Partitions cannot be performed fractionally: one either draws the boundary (generating a new categorical state) or one does not. Hence energy comes in whole-partition units. The work function $W$ is the partition depth the electron must climb (into ionisation) to escape the atomic partition minimum; if the photon's partition distinction is smaller than $W$, the distinction cannot lift the electron out of the minimum, and no emission occurs.

The photoelectric threshold is therefore not a mysterious feature of wave-particle duality. It is a direct statement that partitions have integer count.

\subsection{Compton scattering (Compton 1923)}

Compton \citep{compton1923} showed that X-rays scattering off electrons shift their wavelength by
\[
\Delta\lambda = \lambda' - \lambda = \frac{h}{m_e c}(1 - \cos\theta).
\]

\paragraph{Framework interpretation.} The photon (a partition-operation carrier) and the electron (a bounded oscillatory mode) exchange partition propagation content at rate $c$ when they interact. The wavelength shift formula follows from relativistic partition-state balance: the photon's propagation momentum transfers partially to the electron, with $h/m_e c = 2.426\times 10^{-12}$~m the Compton wavelength setting the scale. The formula is identical to the standard derivation; what differs is interpretation. What is exchanged is not ``wave amplitude'' or ``particle momentum'' separately but a unified partition-transfer content.

\section{Electron Diffraction and the Dissolution of Wave--Particle Duality}
\label{sec:diffraction}

Davisson and Germer \citep{davissongermer1927} scattered electrons off a nickel crystal and observed diffraction peaks at the Bragg angles predicted by $\lambda = h/p$ --- de Broglie's relation \citep{deBroglie1924}. Electrons, seemingly particles, exhibit wave-like diffraction.

\paragraph{Framework interpretation.} The ``wave'' is the non-actualisation structure. An electron at point $A$ is by necessity ``not at points $B, C, D, \dots$'' (Proposition~\ref{prop:arrow}). The spatial distribution of those not-at positions is a field over space: the non-actualisation density. In periodic media the non-actualisation density has periodic structure (the periodicity inherited from the lattice). Interference between overlapping non-actualisation structures from multiple crystal planes produces the observed diffraction pattern.

In this reading, wave--particle duality is not a duality. There is one thing: a bounded oscillatory system (the electron), whose identity at any instant includes its ``not-at'' structure. The ``not-at'' structure has wave properties; the electron itself is always somewhere specific. Measuring position asks ``where is it?''; measuring diffraction asks ``what is its not-at structure?''. These are complementary questions about the same object, and no ontological duality is required.

\section{Quantisation Experiments}
\label{sec:quantisation}

\subsection{Stern--Gerlach (1922)}

Silver atoms passed through an inhomogeneous magnetic field split into exactly two beams \citep{sterngerlach1922}. Not a continuous distribution, not three beams, but two.

\paragraph{Framework interpretation.} Chirality $s$ is a $\mathbb{Z}_2$ topological invariant arising from $\pi_1(\mathrm{SO}(3)) = \mathbb{Z}_2$. Only two values exist: $s = \pm 1/2$. Binary outcomes are \emph{forced}, not probabilistic. The Stern--Gerlach observation confirms that chirality partition is a two-valued topological feature, exactly as required by the partition coordinates of Theorem~\ref{thm:partcoords}.

\subsection{Zeeman effect (1897)}

Zeeman \citep{zeeman1897} observed spectral lines splitting in a magnetic field. A line at angular momentum $l$ splits into $2l+1$ components.

\paragraph{Framework interpretation.} Orientation $m \in \{-l, -l+1, \dots, +l\}$ labels orientations of the partition along the field direction. The SO(3) representation of dimension $2l+1$ is the discrete count of independent orientations at angular complexity $l$. The Zeeman multiplicity is therefore the direct count of partition orientations.

\subsection{Franck--Hertz (1914)}

Franck and Hertz \citep{franckhertz1914} bombarded mercury vapour with electrons and observed that electrons lose energy in discrete quanta of $4.9$~eV.

\paragraph{Framework interpretation.} The mercury atom's partition transitions are discrete. An incoming electron can transfer energy only in amounts equal to a partition state change. There are no continuous states to absorb continuous energy. The $4.9$~eV quantum is the partition depth difference between the Hg ground state and its first allowed excited partition state. Below this threshold, the incoming electron's energy is reflected elastically; at and above, inelastic partition-transition absorption becomes possible. The ``staircase'' structure of the Franck--Hertz current versus voltage curve is the spatial record of integer partition transitions along the electron's path.

\section{Sub-Nucleon Experiments (Without Naming Constituents)}
\label{sec:dis}

Deep inelastic electron--nucleon scattering \citep{friedman1972} at very high momentum transfer $Q^2$ reveals point-like scattering centres within the nucleon. The scaling of structure functions suggests free point-like scatterers at high resolution. Yet no free sub-nucleon particle is ever observed in an asymptotic state.

\paragraph{Framework interpretation.} At sufficient $Q^2$, the probe resolves the sub-coordinate virtual substates transiently within the nucleon's global state. Point-like scattering is expected: the sub-coordinates act as localised (though non-physical) scattering centres within the nucleon. Asymptotic freedom of these sub-coordinates is a framework prediction: as $Q^2$ increases, the probe enters a regime where the sub-coordinates are momentarily behaving as if independent.

Confinement --- the non-observation of free sub-nucleon asymptotic states --- is not an empirical puzzle. It is a theorem: sub-coordinates generically lie outside $[0,1]^3$ and are structurally unobservable in isolation (Theorem~\ref{thm:virt}, Proposition~\ref{prop:subnucvirt}). The framework predicts that no experiment, at any energy, will ever detect a free sub-nucleon particle, because such a detection would require the particle to exit $[0,1]^3$.

% ============================================================
%                        PART IX
% ============================================================
\part{Discussion and Conclusion}

\section{Discussion}
\label{sec:disc}

\subsection{Summary}

We have derived the entire atom from a single axiom.

\begin{itemize}[leftmargin=2em,itemsep=2pt]
\item \textbf{Sub-nucleon scale.} The nucleon's global state admits a recursive sub-coordinate decomposition producing virtual substates that generically lie outside the physical manifold $[0,1]^3$ and are therefore structurally unobservable in isolation. This reproduces ``confinement'' without invoking a colour gauge theory, without introducing new fundamental constants, and without naming sub-particles.
\item \textbf{Nucleon.} A bounded oscillatory system of radius $\approx 0.84$~fm with exponential internal charge distribution, dipole form factor $\Lambda^2 \approx 0.71$~GeV$^2$, matching Hofstadter's measurements.
\item \textbf{Nucleus.} A close-packed bound state of nucleons, with $R = r_0 A^{1/3}$, $r_0\approx 1.2$~fm, binding energies derived from the Composition Theorem reproducing the AME2020 table, and magic numbers $2, 8, 20, 28, 50, 82, 126$ from the nuclear partition shell structure with geometric spin--orbit coupling.
\item \textbf{Charge.} A boundary property emerging only when a partition is drawn. Individual nucleons carry no charge internally; only the nuclear partition boundary carries $+Ze$. This dissolves the intra-nuclear Coulomb repulsion paradox and eliminates the need for a strong force.
\item \textbf{Electron.} A bounded oscillatory mode at the partition depth minimum around the nuclear boundary. The effective $1/r$ potential is the SO(3) partition depth landscape; no Coulomb force is invoked.
\item \textbf{Atom.} Electronic shells derived from capacity $C(n) = 2n^2$; Rydberg formula, aufbau, selection rules, Slater screening all follow.
\item \textbf{Atomic volume.} Most of it is non-actualisation structure --- the electron's ``not-at'' positions. The atom is emptier than matter; it is structured absence.
\end{itemize}

\subsection{Parsimony relative to the standard account}

The standard account of the atom requires, independently:
\begin{itemize}[leftmargin=2em,itemsep=2pt]
\item The electromagnetic interaction (Coulomb's law).
\item The strong interaction (colour $\mathrm{SU}(3)$ QCD).
\item The weak interaction (beta decay, flavour changing).
\item Wave--particle duality as a basic postulate.
\item The Pauli exclusion principle as a basic postulate.
\item The existence of charge $e$ as a parameter.
\item A confinement mechanism (asymptotic freedom + infrared slavery) requiring non-perturbative colour dynamics.
\end{itemize}

The framework requires:
\begin{itemize}[leftmargin=2em,itemsep=2pt]
\item The bounded phase space axiom (Axiom~\ref{ax:bps}).
\item The Loschmidt principle (relational identity, \S\ref{sec:loschmidt}).
\end{itemize}

The strong, electromagnetic, and weak forces; wave--particle duality; exclusion; charge quantisation; and confinement all emerge as theorems. Predictively the two frameworks agree within experimental precision wherever both apply; the partition framework additionally provides a derivation for phenomena (confinement, charge emergence, wave-particle duality) that are postulates or empirical regularities in the standard account.

\subsection{Limitations}

\begin{itemize}[leftmargin=2em,itemsep=2pt]
\item Fundamental particle mass values (electron, proton, neutron masses) are not computed here. In the framework they are fixed points of a non-linear negation operator; their numerical values require solving a fixed-point problem that is algorithmically well-defined but not addressed in this paper. We take measured masses as input.
\item Fine-structure splittings beyond the gross Bohr levels have been compressed: the relativistic and spin--orbit corrections are derivable in the framework via the Partition Lagrangian but are not expanded here for space.
\item Electroweak unification is addressed in companion work; we have focused on the atom here.
\item The inclusion of gravitation and cosmological scales is beyond the scope of the present derivation, although the framework extends in that direction.
\end{itemize}

\subsection{Falsifiable predictions}

The framework is falsifiable and has specific predictions:

\begin{enumerate}[leftmargin=2em,itemsep=2pt]
\item Any experimental outcome distinguishing a \emph{force} from a \emph{gradient of partition depth} would falsify the framework. None has been found in a century of searching.
\item Any isolated sub-nucleon particle --- observed as a free, asymptotic final state without hadronisation --- would falsify the framework. None has been observed in any high-energy experiment to date.
\item Any observation of charge on an unpartitioned entity (e.g.\ a charge measurement on an object with no boundary) would falsify the framework. No such observation exists.
\item Any continuous spectrum from a bounded system (as opposed to a quasi-continuous spectrum from a large-$N$ system) would falsify the framework. No such spectrum has been observed.
\item Any non-integer charge in an isolated particle would falsify the framework. None has been observed.
\item Any deviation from $R = r_0 A^{1/3}$ or from the Hofstadter dipole form factor at the nucleon scale beyond boundary-profile corrections would demand revision. Observed deviations are consistent with the smooth partition boundary profile.
\end{enumerate}

\section{Conclusion}
\label{sec:conc}

The complete atom --- from sub-nucleon virtual substates through the nucleon, the nucleus, the electron, and the observable atomic volume --- derives from a single axiom. No forces are postulated. No strong interaction is needed to hold the nucleus together because no intra-nuclear Coulomb repulsion exists to be overcome: inside the nucleus, the nucleons are unpartitioned and therefore uncharged. Charge emerges only at partition boundaries, which is why Moseley's law depends on $Z$ alone and why fission creates the charges of its fragments. Wave--particle duality dissolves into the single structural observation that an electron's identity includes both its position and its ``not-at'' positions.

Rutherford's $\alpha$ particles showed us that the atom is mostly empty. We now understand \emph{why}: most of the atomic volume is where the electron is not, and it is this non-actualised region that anchors the electron's identity. The emptiness is not the absence of something; it is the structure of what the electron is not.

Mass at every scale --- nucleon, nucleus, atom --- is the accumulated record of what did not happen during the system's existence. Binding energies are partition depth deficits released as partition sharing replaces independence. Magic numbers are the closed-shell occupations of the nuclear partition with geometric spin--orbit coupling. Atomic spectra are the discrete partition transitions of the electron's oscillatory mode around the nuclear boundary.

The framework carries zero adjustable parameters for its structural predictions. The measured quantities $m_p$, $m_n$, $m_e$, $e$, $\hbar$, $c$ are inputs; the radii, binding energies, charge distributions, shell structures, and cross sections are outputs. A century of foundational experiments --- Rutherford (1911), Moseley (1913), Millikan (1913), Einstein photoelectric (1905), Compton (1923), Franck--Hertz (1914), Zeeman (1897), Stern--Gerlach (1922), Davisson--Germer (1927), Hofstadter (1950s--60s), nuclear binding data, magic number observations, deep inelastic scattering --- does not merely fail to contradict the framework but supplies its direct evidence.

The problem was never to hold the nucleus together against its own charge. The problem was that the charge inside the nucleus was never there. Once the geometry is correct, the atom assembles itself.

% ============================================================
%  References
% ============================================================
\bibliography{references}

\end{document}
